% PAGES: 232-234
% Continues by proving the Lemma 7.7 stated at the end of sec07_c3_3.tex
% (PAGES 230-231). Opens directly with "\begin{proof}" -- do not re-open the
% \begin{lemma} for Lemma 7.7, which already closed in the previous file.

\begin{proof}
Since $\bfX \sim N(\mu_{\bfx},\Sigma_{\bfx})$, we need to show that
\begin{equation}
\mu_{\bfx} + U^t\bfZ \;\sim\; N(\mu_{\bfx},\Sigma_{\bfx}),
\end{equation}
in order to conclude that $\bfX = \mu_{\bfx} + U^t\bfZ$.

From Definition 7.4, it follows that $\mu_{\bfx} + U^t\bfZ$ is a multivariate normal random
variable. Recall from Definition 6.1 that, if $U$ is the Cholesky factor of the symmetric
positive definite matrix $\Sigma_{\bfx}$, then
\begin{equation}
\Sigma_{\bfx} \;=\; U^tU.
\end{equation}
Since $\mu_{\bfz}=\bfzero$ and $\Sigma_{\bfz}=I$, it follows from (7.84) and (7.115) that
\begin{align*}
\E{\mu_{\bfx}+U^t\bfZ} &= \mu_{\bfx}; \\
\Sigma_{\mu_{\bfx}+U^t\bfZ} &= U^t\Sigma_{\bfz}(U^t)^t \;=\; U^tU \;=\; \Sigma_{\bfx},
\end{align*}
which is what we wanted to show; cf. (7.114).
\end{proof}

\medskip
\noindent\textit{Example:}\ A two dimensional multivariate normal random variable is called a
\textit{bivariate normal random variable}. We use formula (7.113) to find the probability
density function of a nondegenerate bivariate normal random variable; see (7.118).

Let $\bfX = \begin{pmatrix} X_1 \\ X_2 \end{pmatrix}$ be a nondegenerate bivariate normal
variable, where $X_1$ and $X_2$ are normal random variable with correlation $\rho$ with $\rho
\neq 1$ and $\rho \neq -1$, and with mean and standard deviation $\mu_1, \sigma_1$, and $\mu_2,
\sigma_2$, respectively. Recall from (7.12) that the covariance matrix of $X_1$ and $X_2$ is
\[
\Sigma_{\bfx} \;=\; \begin{pmatrix} \sigma_1^2 & \sigma_1\sigma_2\rho \\ \sigma_1\sigma_2\rho & \sigma_2^2 \end{pmatrix}.
\]
Then,
\begin{equation}
\det(\Sigma_{\bfx}) \;=\; \sigma_1^2\sigma_2^2(1-\rho^2),
\end{equation}
see (10.6), and, since $\rho \neq 1$ and $\rho \neq -1$, we find from (10.7) that
\[
(\Sigma_{\bfx})^{-1} \;=\; \frac{1}{\sigma_1^2\sigma_2^2(1-\rho^2)}
\begin{pmatrix} \sigma_2^2 & -\sigma_1\sigma_2\rho \\ -\sigma_1\sigma_2\rho & \sigma_1^2 \end{pmatrix}.
\]
Let $x-\mu_{\bfx} = \begin{pmatrix} x_1-\mu_1 \\ x_2-\mu_2 \end{pmatrix}$. From (10.20), it
follows that
\begin{align}
(x-\mu_{\bfx})^t(\Sigma_{\bfx})^{-1}(x-\mu_{\bfx})
&= \frac{\sigma_2^2(x_1-\mu_1)^2 - 2\sigma_1\sigma_2\rho(x_1-\mu_1)(x_2-\mu_2) + \sigma_1^2(x_2-\mu_2)^2}{\sigma_1^2\sigma_2^2(1-\rho^2)} \notag \\
% NOTE: source prints a raised "-1" as the numerator of this fraction (verified
% at 600 dpi against PDF page 232, folio 216), which conflicts with the plain
% "-1/2" already applied to this same quantity in the exponent of (7.118) --
% the two negative signs together would make the exponent positive. Reproduced
% exactly as printed; not corrected.
&= \frac{-1}{1-\rho^2}\left(\frac{(x_1-\mu_1)^2}{\sigma_1^2} - \frac{2\rho(x_1-\mu_1)(x_2-\mu_2)}{\sigma_1\sigma_2} + \frac{(x_2-\mu_2)^2}{\sigma_2^2}\right).
\end{align}

Using (7.116) and (7.117), we obtain that the probability density function of the bivariate
normal random variable given by (7.113) can be written as follows:
\begin{equation}
f(x_1,x_2) \;=\; \frac{1}{2\pi\sigma_1\sigma_2\sqrt{1-\rho^2}}\cdot
\exp\!\left(-\frac{1}{2(1-\rho^2)}\left(\frac{(x_1-\mu_1)^2}{\sigma_1^2} - \frac{2\rho(x_1-\mu_1)(x_2-\mu_2)}{\sigma_1\sigma_2} + \frac{(x_2-\mu_2)^2}{\sigma_2^2}\right)\right)
\end{equation}
% NOTE: source does not print a closing square after this Example (verified at
% 600 dpi: PDF page 232, folio 216, ends here; PDF page 233, folio 217, opens
% directly with "Lemma 7.8." and no stray square at its top). Reproduced as
% printed; no square added.

\begin{lemma}
If every two components of a nondegenerate multivariate normal random variable are
uncorrelated, then all the components of the multivariate normal random variable are
independent.

In other words, if $\bfX = (X_i)_{i=1:n}$ is a multivariate normal random variable and if
$\cov(X_j,X_k) = 0$ for all $1 \le j \neq k \le n$, then $X_1, X_2, \ldots, X_n$ are independent
normal random variables.
\end{lemma}

\begin{proof}
Since $\Sigma_{\bfx}(X_j,X_k) = \cov(X_j,X_k) = 0$ for all $1 \le j \neq k \le n$, it follows
that $\Sigma_{\bfx}$ is a diagonal matrix with $\Sigma_{\bfx}(i,i) = \var(X_i) = \sigma_i^2$ for
all $i = 1:n$. Thus, $\Sigma_{\bfx} = \diag\!\left(\sigma_i^2\right)_{i=1:n}$ and
$\Sigma_{\bfx}^{-1} = \diag\!\left(\frac{1}{\sigma_i^2}\right)_{i=1:n}$; see (1.94). From
(10.25), we obtain that
\begin{equation}
(x-\mu_{\bfx})^t(\Sigma_{\bfx})^{-1}(x-\mu_{\bfx}) \;=\; \sum_{i=1}^n \frac{1}{\sigma_i^2}(x_i-\mu_i)^2.
\end{equation}

Also, from (10.1), we find that
\begin{equation}
\det(\Sigma_{\bfx}) \;=\; \prod_{i=1}^n \sigma_i^2.
\end{equation}

From (7.113), and using (7.119) and (7.120), it follows that the probability density function
of the nondegenerate multivariate normal variable $\bfX$ is
\begin{align*}
f(x) &= \frac{1}{(2\pi)^{n/2}\sqrt{\prod_{i=1}^n\sigma_i^2}}\exp\!\left(-\frac12\sum_{i=1}^n\frac{(x_i-\mu_i)^2}{\sigma_i^2}\right) \\
&= \prod_{i=1}^n \frac{1}{\sqrt{2\pi\sigma_i^2}}\exp\!\left(-\frac{(x_i-\mu_i)^2}{2\sigma_i^2}\right) \\
&= \prod_{i=1}^n f_i(x_i),
\end{align*}
where
\begin{equation}
f_i(x_i) \;=\; \frac{1}{\sqrt{2\pi\sigma_i^2}}\exp\!\left(-\frac{(x_i-\mu_i)^2}{2\sigma_i^2}\right).
\end{equation}

Since $f(x)$ can be written as a product of functions of the underlying variables $x_1, x_2,
\ldots, x_n$, we conclude that the random variables $X_1, X_2, \ldots, X_n$ are independent.
Also, since the function $f_i$ given by (7.121) is the probability density function of a normal
variable with mean $\mu_i$ and variance $\sigma_i^2$, see (7.110), it follows that $X_i$ is a
normal random variable, i.e., $X_i \sim N(\mu_i,\sigma_i^2)$, for all $i = 1:n$.
\end{proof}

\begin{theorem}
A random variable $\bfX = (X_i)_{i=1:n}$ is a multivariate normal random variable if and only if
any linear combination of the components $X_1, X_2, \ldots, X_n$ of $\bfX$ is a normal random
variable, i.e., if and only if $\sum_{i=1}^n c_iX_i$ is normal for any $c_i \in \R$, $i = 1:n$.
\end{theorem}

This result is very important for practical applications, e.g., for modeling portfolio returns,
and can be regarded as an alternative definition for multivariate normal variables. Its proof is
beyond our purpose here and can be found, e.g., in \cite{24}.

\begin{theorem}
Let $\bfX \sim N(\mu_{\bfx},\Sigma_{\bfx})$ be an $n$-dimensional multivariate normal random
variable. Let $\bfY$ be the multivariate random variable given by
\begin{equation}
\bfY \;=\; b + M\bfX,
\end{equation}
where $M$ is an $m \times n$ matrix and $b$ is an $m \times 1$ column vector. Then, $\bfY$ is an
$m$-dimensional multivariate normal variable with mean $b + M\mu_{\bfx}$ and covariance matrix
$M\Sigma_{\bfx}M^t$, i.e.,
\begin{equation}
\bfY \;\sim\; N(b+M\mu_{\bfx}, M\Sigma_{\bfx}M^t).
\end{equation}
\end{theorem}

\begin{proof}
Recall from (7.84) that if $\bfY = b+M\bfX$, then
\[
\mu_{\bfY} \;=\; b + M\mu_{\bfx} \qquad \text{and} \qquad \Sigma_{\bfY} \;=\; M\Sigma_{\bfx}M^t.
\]
Thus, in order to establish (7.123), it is enough to show that $\bfY$ is multivariate normal.
We do so by showing that any linear combination of the components $Y_1, Y_2, \ldots, Y_m$ of
$\bfY$ is a normal random variable; see Theorem 7.7.

Let $\sum_{j=1}^m c_jY_j$ be a linear combination of $Y_1, Y_2, \ldots, Y_m$, where $c_j$, $j =
1:m$, are arbitrary constants. Let $C = (c_j)_{j=1:m}$ be the corresponding $m \times 1$ column
vector. From (1.6), we find that
\begin{equation}
\sum_{j=1}^m c_jY_j \;=\; C^t\bfY,
\end{equation}
and therefore, from (7.122) and (7.124), we obtain that
\begin{align}
\sum_{j=1}^m c_jY_j &= C^t\bfY \;=\; C^t(b+M\bfX) \;=\; C^tb+C^tM\bfX \notag \\
&= C^tb+(M^tC)^t\bfX \notag \\
&= \widetilde{c}+\widetilde{C}^t\bfX \notag \\
&= \widetilde{c}+\sum_{i=1}^n \widetilde{c}_iX_i,
\end{align}
where $\widetilde{c} = C^tb$ and $\widetilde{C} = M^tC = (\widetilde{c}_i)_{i=1:n}$.

% NOTE: source reads "any linear combination of Y1, Y2, . . . , Yn" here
% (verified at 600 dpi against PDF page 234, folio 218), even though Y was
% introduced above as m-dimensional with components Y1,...,Ym. Reproduced
% verbatim.
Since $\bfX$ is a multivariate normal random variable, any linear combination of $X_1, X_2,
\ldots, X_n$ is a normal random variable; cf. Theorem 7.7. Thus, from (7.125), it follows that
any linear combination of $Y_1, Y_2, \ldots, Y_n$ is also a normal random variable, and we
conclude, using Theorem 7.7, that $\bfY$ is a multivariate normal random variable.
\end{proof}

\begin{lemma}
Let $\bfZ \sim N(\bfzero,I)$ be an $n$-dimensional multivariate standard normal variable, and
let $Q$ be an $n \times n$ orthogonal matrix. Then, $Q\bfZ$ is also an $n$-dimensional
multivariate standard normal variable, i.e., $Q\bfZ \sim N(\bfzero,I)$.
\end{lemma}

\begin{proof}
Note that, if $\bfZ \sim N(\bfzero,I)$, then $\mu_{\bfz} = \bfzero$ and $\Sigma_{\bfz} = I$. From
(7.123), it follows that
\[
Q\bfZ \;\sim\; N(\bfzero,Q\Sigma_{\bfz}Q^t) \;=\; N(\bfzero,QQ^t) \;=\; N(\bfzero,I),
\]
since $QQ^t = I$ for any orthogonal matrix $Q$; cf. (10.14).
\end{proof}

% END OF CHUNK c3 (PAGES 227-234). PDF page 234 (folio 218) ends cleanly here,
% immediately after the proof of Lemma 7.9 closes -- no environment and no
% non-environment example/solution block is left open. Section 7.6
% "Multivariate normal random variables" itself continues past this point
% (chunk c4 owns PDF pages 235-242); do not re-open its \section{} heading.
